\chapter{<<vi:Kiến trúc hệ thống||en:System Architecture>>}
\label{ch:arch}
<<vi:Kiến trúc hệ thống mô tả các thành phần và cách chúng trao đổi dữ liệu. Bắt đầu từ đầu vào và đầu ra, sau đó giải thích trách nhiệm của từng thành phần; chỉ nêu công cụ khi công cụ đó giúp làm rõ một lựa chọn thiết kế.||en:The architecture describes the components and their data exchange. Start with inputs and outputs, then explain each component’s responsibility. Mention a tool when it clarifies a design choice.>>

%%IF diagram
\section{<<vi:Tổng quan||en:Overview>>}
\begin{figure}[htbp]
  \centering
  \begin{tikzpicture}[lk]
    \node[layer=11cm]                 (ui)   {<<vi:Giao diện||en:Interface>>};
    \node[layer=11cm, below=4mm of ui] (core) {<<vi:Xử lý lõi||en:Core processing>>};
    \node[layer=11cm, below=4mm of core, lkcolor=lkteal] (data) {<<vi:Dữ liệu||en:Data>>};
  \end{tikzpicture}
  \caption{<<vi:Kiến trúc nhiều tầng (style \texttt{layer})||en:Layered architecture (\texttt{layer} style)>>}
  \label{fig:layers}
\end{figure}

\section{<<vi:Luồng xử lý||en:Processing flow>>}
\begin{figure}[htbp]
  \centering
  \lkpipeline{<<vi:Thu thập||en:Collect>>, <<vi:Làm sạch||en:Clean>>, <<vi:Huấn luyện||en:Train>>, <<vi:Đánh giá||en:Evaluate>>}
  \caption{Pipeline}
\end{figure}

\begin{figure}[htbp]
  \centering
  \lkdiagram{diagrams/architecture}
  \caption{<<vi:Sơ đồ D2||en:D2 diagram>>: \lkfile{diagrams/architecture.d2}}
\end{figure}
%%ENDIF
